\documentclass{uinwsskripsi}
\usepackage[utf8]{inputenc}
\usepackage[T1]{fontenc}
\usepackage{geometry}
\usepackage{graphicx}
\usepackage{float}
\usepackage{hyperref}
\usepackage[english,indonesian]{babel}

\title{Optimasi Pemrosesan Template LaTeX Menggunakan Pendekatan Berbasis AI}
\author{Academic Support Agent}
\date{\today}

\begin{document}
\maketitle

\begin{abstract}

\end{abstract}

\chapter{PENDAHULUAN}


\chapter{TINJAUAN PUSTAKA}


% Original content (commented out):
% \documentclass{uinwsskripsi}

\input{data_skripsi.tex}

\begin{document}

\input{muka_skripsi.tex}
\input{bab-1-skripsi.tex}
\input{bab-2-skripsi.tex}
\input{bab-3-skripsi.tex}
\input{bab-4-skripsi.tex}
\daftarpustaka
\lampiran
\input{lampiran-1.tex}
\input{lampiran-2.tex}

\end{document}

\end{document}